% Laget automatisk av elevark/lag_elevark.py fra sim/lodd-i-snorer.html. Endringer her blir overskrevet.
\documentclass[11pt]{article}
\usepackage{elevark}
\begin{document}
\elevarkhode{Fysikk 2 · Bevegelse og krefter i to dimensjoner}{Lodd i to snorer}{https://lektorbo.no/sim/lodd-i-snorer.html}{Simuleringen}

Et lodd henger i en knute som holdes oppe av to snorer festet i taket. Når loddet henger i ro, er summen av kreftene på knuten null, både vannrett og loddrett. Hvor store blir snorkreftene, og hvordan avhenger de av vinklene?

\section*{Oppgaver}

Start med $α = β = 30^{\circ}$ og 2,0~kg.

\begin{oppgaveliste}
\oppgave \fase{gjett}{Gjett}Loddet veier 19,6~N, og det er to snorer. Er hver snorkraft halvparten av tyngden, mindre eller større?
\svar
\oppgave \fase{test}{Regn}Sett opp likevekt vannrett og loddrett. Vis at $S = G / (2 \sin  α)$ når vinklene er like, og regn ut snorkraften for 30°. Sjekk.
\svar
\oppgave \fase{test}{Undersøk}Gjør begge vinklene mindre og mindre. Hva skjer med snorkreftene? Regn ut snorkraften for 5°.
\svar
\oppgave \fase{forklar}{Forklar}Hvorfor er det umulig å stramme en klessnor helt rett når det henger klær på den?
\svar
\oppgave \fase{test}{Regn}Sett $α = 60^{\circ}$ og $β = 30^{\circ}$. Hvilken snor har størst kraft? Regn ut begge snorkreftene, og sjekk at de vannrette komponentene er like store.
\svar
\oppgave \fase{forklar}{Forklar}Se på trekanten til høyre. Hvorfor må kreftene danne en lukket trekant når knuten er i ro? Hva er vinkelen mellom $S_{1}$ og $S_{2}$ i trekanten?
\svar
\oppgave \fase{utfordring}{Utfordring}En snor tåler 200~N. Hvor liten kan vinkelen være når det henger 5,0~kg i midten av snora, med like vinkler på begge sider?
\svar
\end{oppgaveliste}

\end{document}
